\section*{الفصل \ref{ch:g12:light-as-wave} --- الضوء موجةً: الحيود والتداخل}

\begin{solution}{exo:g12:light-as-wave:1}
$f = c/\lambda$: الأحمر $\num{3.00e8}/\num{7.0e-7} \approx
\qty{4.3e14}{Hz}$؛ والبنفسجي $\approx \qty{7.5e14}{Hz}$ --- فالبنفسجي
أعلى (لأن \omterm{def:g10:light-spectra:wavelength}{طول موجته} أقصر، والسرعة $c$ نفسها).
\end{solution}

\begin{solution}{exo:g12:light-as-wave:2}
$\theta = \num{5.32e-7}/\num{2.0e-4} = \qty{2.7e-3}{rad}$؛
$L = 2\theta D = 2 \times \num{2.66e-3} \times 1.5 \approx
\qty{8.0}{mm}$.
\end{solution}

\begin{solution}{exo:g12:light-as-wave:3}
\omterm{def:g10:signals-and-waves:sound}{الصوت}: $\lambda = 340/425 = \qty{0.80}{m}$، ومنه $\lambda/a = 1.0$ ---
أي \omterm{def:g12:light-as-wave:diffraction}{حيود} أقصى، فيغمر \omterm{def:g10:signals-and-waves:sound}{الصوت} الممرّ. والضوء:
$\lambda/a \approx \num{5.5e-7}/0.80 \approx \num{7e-7}$ ---
أي انتشار مهمَل: فيبقى الضوء مستقيمًا، وتخفيه الزاوية.
\end{solution}

\begin{solution}{exo:g12:light-as-wave:4}
عند $M$: $\delta = \qty{6.0}{cm} = 3\lambda$، \omterm{prop:g12:light-as-wave:conditions}{فتداخل بنّاء} --- أي \omterm{def:g10:signals-and-waves:sound}{صوت} عالٍ.
وعند $N$: $\delta = \qty{5.0}{cm} = (2 + \tfrac12)\lambda$، \omterm{prop:g12:light-as-wave:conditions}{فتداخل هدّام}
--- أي صمت (شبه تامّ).
\end{solution}

\begin{solution}{exo:g12:light-as-wave:5}
$i = \lambda D/a = \num{6.33e-7} \times 1.8/\num{2.5e-4} \approx
\qty{4.6}{mm}$. وضمن $\pm\qty{1.0}{cm}$: $|k| \leq 10/4.6 = 2.2$،
ومنه $k = -2, \dots, 2$: أي خمسة أهداب ساطعة.
\end{solution}

\begin{solution}{exo:g12:light-as-wave:6}
$d = 2\lambda D/L = 2 \times \num{6.5e-7} \times 1.6/0.028 \approx
\qty{74}{\micro m}$.
\end{solution}

\begin{solution}{exo:g12:light-as-wave:7}
$i = \qty{4.7}{mm}$، ومنه $\lambda = i\,a/D = \num{4.7e-3} \times
\num{2.0e-4}/1.5 \approx \qty{6.3e-7}{m} = \qty{630}{nm}$: أي أحمر.
وقياس عشرة أبعاد هدبية يقسم خطأ قراءة المسطرة على عشرة.
\end{solution}

\begin{solution}{exo:g12:light-as-wave:8}
للمصباحين تواتران متوافقان (على وجه العموم) لكنهما \emph{ليسا على
وتيرة واحدة}: فكلٌّ يصدر قطارات موجية قصيرة يقفز طورها عشوائيًّا، ومن ثم
يكون المنبعان غير \omterm{def:g12:light-as-wave:coherent}{مترابطين} وتنزاح الأهداب اللحظية
مليارات المرات في \omterm{def:g10:orders-of-magnitude:unit}{الثانية} --- فترى العين المتوسط، وهو منتظم.
أما يونغ فيضيء الشقّين من \omterm{def:g10:signals-and-waves:wave}{موجة} \emph{واحدة}: فترث النسختان
كل قفزة طور معًا وتبقيان على وتيرة واحدة.
\end{solution}

\begin{solution}{exo:g12:light-as-wave:9}
من الأول إلى السابع $= 6i = \qty{18.0}{mm}$، ومنه $i = \qty{3.0}{mm}$؛
و$\lambda = i\,a/D = \num{3.0e-3} \times \num{3.0e-4}/1.4 \approx
\qty{6.4e-7}{m} \approx \qty{640}{nm}$.
\end{solution}

\begin{solution}{exo:g12:light-as-wave:10}
$\sin\theta = k\lambda/d = 0.396\,k$: $k = 1$، $\theta = 23^\circ$؛
$k = 2$، $\sin\theta = 0.79$، $\theta = 52^\circ$؛ أما $k = 3$ فكان سيقتضي
$\sin\theta = 1.19 > 1$: وهذا مستحيل. فأعلى رتبة $k = 2$.
\end{solution}

\begin{solution}{exo:g12:light-as-wave:11}
$\theta = \num{5.5e-7}/0.10 = \qty{5.5e-6}{rad}$. ويقابل مصباحا السيارة
زاوية $1.5/\num{1.0e5} = \qty{1.5e-5}{rad}$، أي نحو ثلاثة أضعاف
التشويش: فيمكن فصلهما، بالكاد.
\end{solution}

\begin{solution}{exo:g12:light-as-wave:12}
يتداخل انعكاسا الوجهين الأمامي والخلفي؛ وفي \omterm{def:g10:light-spectra:white}{الضوء الأبيض} يعزّز كل
سُمك بعض أطوال \omterm{def:g10:signals-and-waves:wave}{الموجة} ويلغي غيرها، ومن هنا اللون.
وسُمك الغشاء لا يتوقف إلا على الارتفاع، ومن ثم تشكّل النقاط ذات اللون
الواحد أشرطة أفقية. والنزّ يرقّق الغشاء، فيقع السُّمك
الذي رسم لونًا معينًا أخفض فأخفض: فتزحف الأشرطة
إلى أسفل.
\end{solution}

\begin{solution}{exo:g12:light-as-wave:13}
$i = \lambda D/a$: البنفسجي $\num{4.0e-7} \times 1.5/\num{2.0e-4} =
\qty{3.0}{mm}$؛ والأحمر $\qty{5.6}{mm}$. وعند المركز: كل الألوان ساطعة ---
أي هدب أبيض واحد. وعلى بعد بضعة مِلّيمترات تكون أنماط الألوان قد
تباعدت: فأهداب متقزّحة. وأبعد، تتراكب عظميات كل الألوان
في كل موضع: فتشوّش أبيض منتظم.
\end{solution}

\begin{solution}{exo:g12:light-as-wave:14}
$\theta = \num{5.5e-7}/2.4 = \qty{2.3e-7}{rad}$؛ وعند \qty{250}{km}
يكون ذلك $\num{2.3e-7} \times \num{2.5e5} \approx \qty{6}{cm}$. فيقدر أن
يعدّ السيارات ويلتقط شخصًا، أما حروف الصحف (وهي مِلّيمترية) فأدنى
عشرين مرة من حدّ \omterm{def:g12:light-as-wave:diffraction}{الحيود}: فالادّعاء أسطورة.
\end{solution}

\begin{solution}{exo:g12:light-as-wave:15}
\omterm{def:g10:signals-and-waves:frequency}{التواتر} يثبّته المنبع: فلا يتغير. أما السرعة فتهبط إلى
$c/n$، ومنه $\lambda' = \lambda/n$ و
$i' = i/n = 4.0/1.33 \approx \qty{3.0}{mm}$: فتتضايق الأهداب.
\end{solution}

\begin{solution}{pb:g12:light-as-wave:1}
\textbf{1.} $\theta = \lambda/a = \num{6.33e-7}/\num{1.00e-4} =
\qty{6.3e-3}{rad}$.

\textbf{2.} نصف العرض $D\tan\theta \approx D\theta = \lambda D/a$،
ومنه $L = 2\lambda D/a = 2 \times \num{6.33e-7} \times
3.00/\num{1.00e-4} = \qty{3.8}{cm}$.

\textbf{3.} $(3.8 - 3.7)/3.8 \approx 3\%$: أي في حدود دقة
المسطرة --- فالطريقة صحيحة.

\textbf{4.} $L \propto 1/a$: فتنصيف $a$ يضاعف $L$، ومنه
$L = \qty{7.6}{cm}$. فكلما ضاق العائق أو الشقّ اتسع النمط.

\textbf{5.} $\tan(\num{6.33e-3}) = \num{6.33008e-3}$: فهما يتوافقان في حدود
نحو $\num{1e-5}$ نسبيًّا --- أي أن التقريب أفضل بكثير من
أي قياس هنا.

\textbf{6.} بحسب \cref{prop:g12:light-as-wave:hair}: العائق الذي عرضه
$d$ يحيد كشقّ عرضه $d$، ومن ثم يبقى $L = 2\lambda D/d$
صحيحًا.

\textbf{7.} $d = 2\lambda D/L = 2 \times \num{6.33e-7} \times
3.00/0.054 \approx \qty{70}{\micro m}$.

\textbf{8.} $L = \qty{5.2}{cm} \to d = \qty{73}{\micro m}$؛
$L = \qty{5.6}{cm} \to d = \qty{68}{\micro m}$: ومنه
$d = \qty{70 \pm 3}{\micro m}$.

\textbf{9.} $d = \num{3.80e-6}/0.076 = \qty{50}{\micro m}$: فشعرة
الصديق أرقّ --- و$L \propto 1/d$، ومن ثم ترمي الشعرة الأرقّ
النمط الأوسع.

\textbf{10.} لا: فالمصباح أبيض (فيرسم كل طول \omterm{def:g10:signals-and-waves:wave}{موجة} نمطًا مختلفًا،
ويتشوّشان) وغير مترابط وعريض (فلا حزمة واحدة نظيفة
تحيد). أما \omterm{def:g11:color-light-sources:lamps}{الليزر} \omterm{def:g10:light-spectra:white}{فأحادي اللون} وموجَّه.

\textbf{11.} لا يترابط إلا نسختان من \omterm{def:g10:signals-and-waves:wave}{موجة} واحدة؛ فمنبعان مستقلان
ينحرفان عن الوتيرة وتنمحي الأهداب
(\cref{rem:g12:light-as-wave:lamps}).

\textbf{12.} عند المركز $\delta = 0 = 0 \times \lambda$ مهما يكن
$\lambda$: \omterm{prop:g12:light-as-wave:conditions}{فالتداخل بنّاء} مهما يكن \omterm{def:g12:mechanical-waves:wavelength}{طول الموجة}.

\textbf{13.} $i = 4.26/10 = \qty{4.3}{mm}$؛ فقياس امتداد عشرة يُقرأ
بخطأ المسطرة نفسه الذي يُقرأ به بعد هدبي واحد، ومن ثم يُقسم الخطأ على $i$
على عشرة.

\textbf{14.} $\lambda = i\,a'/D = \num{4.26e-3} \times
\num{2.50e-4}/2.00 = \num{5.33e-7} \approx \qty{533}{nm}$.

\textbf{15.} يقع $\qty{533}{nm}$ داخل $\qty{532 \pm 10}{nm}$:
فهو متسق؛ وذلك الطول أخضر.

\textbf{16.} $\tan\theta = 0.43/1.00$، ومنه $\theta = 23^\circ$. وهنا
$\tan\theta = 0.43$ بينما $\sin\theta = 0.40$: فعند $23^\circ$ يخطئ
التماهي في الزوايا الصغيرة بين $\sin$ و$\tan$ و$\theta$ بنحو
$10\%$ --- فاستعمل الدوال الدقيقة.

\textbf{17.} $d = \lambda/\sin\theta = \num{6.33e-7}/0.396 \approx
\qty{1.6}{\micro m}$.

\textbf{18.} $\tan\theta = 1.65$، $\theta = 59^\circ$،
$\sin\theta = 0.855$: $d = \num{6.33e-7}/0.855 \approx
\qty{0.74}{\micro m}$.

\textbf{19.} $1.6/0.74 \approx 2.2$: فمسارات قرص DVD أكثف مرتين،
وحفره أصغر بالمقابل (وتُقرأ \omterm{def:g11:color-light-sources:lamps}{بليزر} أقصر طول \omterm{def:g10:signals-and-waves:wave}{موجة}) ---
أي أضعاف المعطيات في القرص \qty{12}{cm} نفسه.

\textbf{20.} الشعرة \qty{70 \pm 3}{\micro m}؛ وشعرة الصديق
\qty{50}{\micro m}؛ \omterm{def:g11:color-light-sources:lamps}{والليزر} الأخضر \qty{533}{nm}؛ ومسارات CD
\qty{1.6}{\micro m}؛ ومسارات DVD \qty{0.74}{\micro m} --- فطول \omterm{def:g10:signals-and-waves:wave}{موجة}
معلوم يحوّل هندسة كل نمط إلى طول، فيصير ضوء نصف
الميكرومتر مسطرةً مدرّجة بنصف ميكرومتر.
\end{solution}
